\begin{acknowledgments}
The author used AI assistants in this work. Anthropic Claude was used in two separate sessions. The first session (claude-opus-5-5) planned the study, wrote the instructions for each stage, re-ran every verification suite and checked the results independently, and reviewed the manuscript. The second session (claude-opus-5-5, with Anthropic Claude Fable, claude-fable-5-1 and claude-fable-5-2, for parts of versions 1.1 and 1.2) derived the lemmas, wrote the computer-assisted certificates and the verification code, and drafted the text; it also found and reported weaknesses in its own and the first session's reasoning, which were corrected in versions 1.1 to 1.3. Google Gemini (Gemini 4 Argon) and OpenAI ChatGPT (GPT-5.6 Luna; GPT-6 Sol for version 1.3) independently reviewed the manuscript and re-derived its main steps; their comments led to the line-by-line proofs of the appendices, the specification of the computer-assisted proofs, the lower end of Lemma BE, and several corrections. The author directed the research, made all decisions on its content, and takes full responsibility for this article.
\end{acknowledgments}

\emph{Data and code availability.}
All code, stored results, the computer-assisted certificates and the verification suite are available at \url{https://github.com/neginyan/resolution-gap-weak-ripple} and archived at Zenodo, \href{https://doi.org/10.5281/zenodo.23932577}{doi:10.5281/zenodo.23932577}.
A single entry point, \code{verification/run\_all.py}, re-derives every number of this paper from the stored results and recomputes a sample of every certificate with the rounding-covering acceptance rule of App.~\ref{app:cells} (about twenty minutes on one CPU core); \code{--recompute} recomputes all of them.
Paper~F and its repository~\cite{F} are not changed.
